\documentclass[10pt,a4paper]{article}
\usepackage[margin=2.2cm]{geometry}
\usepackage{amsmath,amssymb}
\setlength{\parskip}{3pt}\setlength{\parindent}{0pt}
\title{Radiation cancellation in a $\pi$ shift sidewall Bragg cavity\\[2pt]
\large Analytical model and status, version 7}
\author{For Evyatar Rubin, Rosenthal group, Technion}
\date{5 October 2026. Replaces v6.2 (27 September) and the review note of 5 October.}
\begin{document}
\maketitle
Labels: \textbf{M} measured, \textbf{D} derived from measured values, \textbf{E} expected. Devices, all $N=98$
per side, mode width 19 to 20\,$\mu$m: A = TE plain (corr 250), B = TE overshoot apodization, C = TM plain
(corr 325). Corrected far fields: \texttt{results\_from\_igum/farfield\_sph\_20um/} (IGUM 100034).

\section{Bottom line}
\begin{enumerate}\setlength{\itemsep}{1pt}
\item \textbf{One formula.} The far field is that of the polarization current $\Delta\varepsilon\mathbf E$
radiating in uniform oxide, Eq.~(\ref{eq:A}). Harmonics, features and cancellation follow by linear algebra.
\item \textbf{Cancelling one spherical harmonic does not work here.} The largest harmonic group of A holds
19\,\%, so at most $1.24\times$ in $Q_{\rm rad}$. The table also depends on the choice of origin (a 0.19\,$\mu$m
shift moves the magnetic dipole from 3 to 19\,\%), because the grating has no mirror symmetry in $x$.
\item \textbf{The right object is the response matrix, not a basis.} Cancellation is a least squares problem
on the whole sphere, Eq.~(\ref{eq:ls}). A dominant ``single term'' exists only if one controllable pattern
(a singular vector of the response) carries most of the leak.
\item \textbf{Plain TE leaks from the junction.} 83\,\% of A's pattern is reproduced by sources within
$\pm0.5$\,$\mu$m of the $\pi$ shift (B 19\,\%, C 42\,\%, null 0). This is a ceiling for redesigning the innermost
teeth, not a prediction; an independent estimate (GPT) of the realisable part is of order 10\,\%.
\item \textbf{The apodized device B is a distributed radiator}: a job for the inverse design, not for a few
knobs. \textbf{Plain TM} sends two thirds of its leak into a grazing needle: a job for a periodic comb.
\item \textbf{Envelope shaping stays the large lever}: B is at least $15\times$ better than A in $Q_{\rm rad}$.
Mode width alone sets no lower bound on the leak.
\item \textbf{Phase.} A feature inside the core has only a sign and a size. Only a comb offset or a position in
the cladding gives a continuous phase. A pure second harmonic of the corrugation puts nothing in the light cone.
\item \textbf{The far field error is fixed.} Flat monitor projections are right only within $40^\circ$ of the
monitor normal. The new projection uses one surface around the waveguide and certifies itself. TE results
changed by less than one point; TM is still uncertain at grazing angles.
\item \textbf{New:} 1.2 to 1.9 points of the port loss do not come from the cavity, most likely end scattering
(E). If so, $Q_i$ does not depend on length and B radiates four times less than its port loss says.
\item \textbf{Lengthening to the Q3dB device should not change the harmonic percentages} (E, one run needed).
Removing a fraction $f$ of the leak multiplies $Q_{3\rm dB}$ by $1/(1-f)$.
\end{enumerate}

\section{Model}
\textbf{Setting.} Axes: $x$ along the guide, $y$ transverse, $z$ vertical; $e^{-i\omega t}$.
$n_{\rm SiN}=1.97$, $n_c=1.444$, uniform oxide outside. At 1559\,nm $k=5.82$\,$\mu$m$^{-1}$; $\beta=6.28$ (TE) and
6.08 (TM), just outside the light cone. $\Delta\varepsilon=1.80$ in SiN; a feature changes it by
$\delta\varepsilon_P=\mp1.80$.

\textbf{Exact far field} (contrast moved to the right of Maxwell's equations, oxide Green tensor):
\begin{equation}\label{eq:A}
\mathbf A(\hat{\mathbf r})=(\mathsf I-\hat{\mathbf r}\hat{\mathbf r})\!\int\!\Delta\varepsilon\,\mathbf E\,
e^{-ik\hat{\mathbf r}\cdot\mathbf r}d^3r,\qquad
\frac{dP}{d\Omega}=\frac{k_0^4}{32\pi^2\eta}|\mathbf A|^2,\qquad Q_{\rm rad}=\frac{\omega U}{P}.
\end{equation}
\textbf{Harmonics.} $\mathbf A=\sum[\alpha_{lm}\mathbf X_{lm}+\beta_{lm}\hat{\mathbf r}\times\mathbf X_{lm}]$,
powers add. In terms of the near field,
\[
\alpha_{lm}=4\pi(-i)^l\!\int\!\Delta\varepsilon\,j_l(kr)\mathbf X^*_{lm}\!\cdot\mathbf E\,d^3r,\qquad
\beta_{lm}=\frac{4\pi i(-i)^l}{k}\!\int\!\Delta\varepsilon\,\mathbf E\cdot\nabla\times[j_l(kr)\mathbf X^*_{lm}]d^3r .
\]
Since $j_l\sim(kr)^l$, order $l$ lives at $r\sim l/k$ ($l=3$ is the first tooth). Selection rules come only from
the $y$ and $z$ mirrors (TE: electric $l+m$ even, magnetic $l+m$ odd; TM the converse). There is no $x$ rule:
both arms are built in the same order, and the standing wave has a node at $x=-0.125$ and an antinode at
$+0.125$\,$\mu$m (M). \emph{This replaces the symmetry section of v6.2.}

\textbf{Adding a feature} (exact to first order):
\begin{equation}\label{eq:split}
\mathbf A_{\rm new}-\mathbf A_c=\underbrace{\mathcal R[\delta\varepsilon_P\mathbf E_0]}_{\text{frozen field}}
+\underbrace{\mathcal R[\Delta\varepsilon\,\delta\mathbf E]}_{\text{redistribution}}
+\delta k\,\partial_k\mathcal R[\Delta\varepsilon\mathbf E_0].
\end{equation}
The first term is the feature radiating in the unchanged field (closed form, for screening). The second is the
change of the cavity field; only a run measures it. A moved sidewall uses $\mathbf E_\parallel$ and $D_\perp$.
\begin{itemize}\setlength{\itemsep}{0pt}
\item \emph{Small feature} ($R\lesssim0.15$\,$\mu$m): a dipole
$\mathbf p=\varepsilon_0\delta\varepsilon_PV_PL\,\mathbf E_0(\mathbf r_p)$ with phase
$e^{-ik\hat{\mathbf r}\cdot\mathbf r_p}$; $L=1.30$ for an oxide hole in SiN. At the origin it is a pure dipole;
displaced, its power spreads over the orders (share in $l=1,2,3,4$): $kx_p=1$: 0.86, 0.13, 0.01, 0;
$kx_p=2$: 0.55, 0.36, 0.08, 0.01; $kx_p=4$: 0.14, 0.29, 0.33, 0.18. One tooth is $kx_p=1.5$, so no feature
radiates a single order.
\item \emph{Pair} at $\pm x_p$: factor $2\cos(kx_p\sin\theta\cos\varphi)$; a sign, no phase.
\item \emph{Comb} of period $\Lambda_2$ and offset $\delta x$: a narrow band at $u_x=(\beta-K_2)/k$ with phase
$\phi_2=2\pi\delta x/\Lambda_2$; to aim at $u_x$ use $\Lambda_2=2\pi/(\beta+k|u_x|)$.
\end{itemize}

\textbf{Cancellation.} With real knobs $\mathbf q$, response columns $\mathsf B=\partial\mathbf A/\partial\mathbf q$,
and $\mathsf N$ the null space of the width and resonance sensitivities:
\begin{equation}\label{eq:ls}
\min_{\mathbf z}\|\mathbf A_c+\mathsf{BN}\mathbf z\|^2,\qquad
\eta^2=\frac{\|\mathsf P_{\mathsf{BN}}\mathbf A_c\|^2}{\|\mathbf A_c\|^2},\qquad
\text{one knob: }\ \frac{P(q)}{P_c}=1-\eta^2+\eta^2\Bigl(\frac q{q_*}-1\Bigr)^2 .
\end{equation}
$\eta^2$ is the removable fraction. A change of basis cannot raise it. The useful ``one term'' is a left singular
vector of $\mathsf{BN}$ that carries most of $\mathbf A_c$; its right singular vector is the list of tooth
changes. Cancelling $K$ complex channels at fixed width and resonance needs $2K+2$ real knobs. Removing a channel
that holds a fraction $f$ gives at most $1/(1-f)$. It gets worse than the baseline for $q>2q_*$ or the wrong
sign, which is likely when the feature feeds other channels, when the target is small, or when restoring the
width rewrites $\mathbf A_c$. The objective must be $P_{\rm rad}/U$.

\textbf{Mode width.} Decay rate $\gamma=\sqrt{\kappa^2-\delta_\beta^2}$, intensity width $\ln2/\gamma$. A feature
shifts the resonance (second order effect on the width at gap centre); features repeating with the pitch change
$\kappa$ at first order (oxide holes in narrow segments shorten the mode, in wide segments lengthen it);
cladding posts at 1.5 to 2\,$\mu$m do neither. A centre feature is compensated most cheaply by the gap length.

\textbf{Envelope.} The source carriers sit at $\pm\beta$, 0.46\,$\mu$m$^{-1}$ outside the cone; the leak is the
envelope's Fourier content beyond that offset plus the junction. For the same 20\,$\mu$m width the scalar cone
content is $9\times10^{-5}$ for $e^{-\gamma|x|}$, $3\times10^{-15}$ for a Gaussian, zero for $\mathrm{sinc}^2$ (D).
The limits are practical: a Gaussian needs $\kappa=0.069$\,$\mu$m$^{-1}$ at the half width (plain: 0.034 to
0.044), plus joins and finite length.

\section{Far field measurement}
\textbf{Error.} The harmonic tables were computed on two flat monitor projections stitched at $45^\circ$. They
disagree there (side/top power 1.3 for TE, 2.3 for TM). Cause, shown with an exact dipole and no simulation: with
monitors 2.15\,$\mu$m from the axis and $\pm3.4$\,$\mu$m wide, a flat projection keeps only 0.73 of the power at
40 to $50^\circ$ from its normal and 0.38 at 50 to $60^\circ$. A larger box is worse.

\textbf{Fix.} $\mathbf E$ and $\mathbf H$ on the crossing parts of the two monitors, with their mirror images,
form one surface around the guide, open at the $x$ ends with a 15\,$\mu$m taper. One coherent surface integral
gives the far field (\texttt{python\_tools/farfield\_surface.py}, engine flag \texttt{save\_surface\_eh}). It
reproduces exact dipoles to 0.2\,\%; GPT reviewed it and one bug it found is fixed. Certificate on every result:
far field power equals the flux through the surface.

\section{Results}
\begin{center}\small\begin{tabular}{lccc}\hline
(M, corrected far field) & A (TE plain) & B (TE oversh.) & C (TM plain)\\\hline
$T$ / $Q_L$ / width ($\mu$m) & 0.912 / 1694 / 19.1 & 0.973 / 7696 / 19.6 & 0.915 / 1652 / 19.2\\
far field power / surface flux & 1.000 & 0.986 & 0.974\\
power through the surface (\% of input) & 7.33 & 0.69 & 6.5\\
port loss $1-T-R$ (\%) & 8.55 & 2.63 & 8.21\\
largest harmonic groups (\%) & E(3,$\pm$3) 19, M(1,0) 14, & M(2,$\pm$1) 6, & M(2,$\pm$2) 14, E(3,0) 9,\\
 & M(2,$\pm$1) 13, E(4,$\pm$4) 12 & M(1,0) 4 & E(1,0) 8\\
share in $l\le5$ / mean $l$ & 95.5 / 3.1 & 31 / 13.8 & 74 / 5.4\\
share at $|u_x|>0.9$ (\%) & 9 & 11 & 67\\
reachable from $\pm0.5$ / $\pm10$\,$\mu$m (D) & 0.83 / 0.88 & 0.19 / 0.74 & 0.42 / 0.50\\
corrected vs stitched (correlation) & 0.978 & 0.970 & 0.932\\\hline
\end{tabular}\end{center}
\begin{itemize}\setlength{\itemsep}{1pt}
\item \textbf{A} is a compact, low order radiator, spread over five groups. Cross check of the junction fit
(fit on one part of the sphere, predict another): $+0.75$, $+0.66$, and one failure ($-0.41$). Supported, not
proven. Which current component radiates is not established.
\item \textbf{B} is high order and distributed (sources at $\pm6$, $\pm13$, $\pm25$\,$\mu$m).
\item \textbf{C} depends on the end taper (6.6 to 5.9\,\%): needs longer monitors before it is quoted.
\item \textbf{End loss.} The gap between rows 3 and 4 is 1.2 to 1.9 points in all three. With the cavity part
alone, $Q_i$(A, $N=98$) $=4.35\times10^4$, equal to $4.41\times10^4$ at $N=166$ (D); for B,
$Q_i\approx2\times10^6$ instead of $6\times10^5$ (E). Not yet tested directly.
\item \textbf{Earlier results that stand} (from v6.2). Comb period: emission predicted at $u_x=-0.989$ for
531\,nm; measured 531 locked ($+16\,\%$ $Q$ at fixed width), 545 harmful, 551 null. Comb phase:
$T$ versus offset is a pure first harmonic, lock at $270^\circ$. TE post pair near the centre: loss 0.119 to
0.093 at fixed width, a uniform scaling of the pattern, that is the redistribution term. Hole lattices change
the width with the predicted sign (10.2 to 25.6\,$\mu$m against a 15.5\,$\mu$m control).
\item \textbf{Superseded} (from v6.2): its TE harmonic content (8\,\% dipole, 28\,\% in $l=3$) and its ceiling of
about 10\,\%, which holds only for on axis features (2\,\% here).
\end{itemize}

\section{What each knob can do (E unless marked)}
\begin{center}\small\begin{tabular}{p{4.2cm}p{11.2cm}}\hline
Innermost teeth, cavity length & The knobs for A. Ceiling 83\,\% (D), realisable part unknown. Strong width and
resonance side effects, removed through $\mathsf N$.\\
On axis holes or posts & 2 to 5\,\% (D). Dead end.\\
Second set of teeth ($2G$) & Carriers at $\pm3\beta$, outside the cone. No direct lever.\\
Tooth shifts along the arms & Continuous phase. The knobs for B; the inverse design uses them.\\
Cladding comb & Narrow band with a phase dial. Suits the TM needle, not the broad TE hump.\\\hline
\end{tabular}\end{center}

\section{Q3dB devices}
At $-3$\,dB, $Q_{3\rm dB}=(1-1/\sqrt2)Q_i=0.2929\,Q_i$; length only sets the coupling.
\begin{itemize}\setlength{\itemsep}{1pt}
\item \textbf{Harmonic percentages when lengthened: predicted unchanged} within 1 to 2 points for the plain
devices (E). The pattern is set by the junction; added mirror periods change the coupling, not the junction
field. The TE Q3dB device is A lengthened from 98 to 166 periods. The radiated share of the input rises from
7.3 to about 41\,\%: same pattern, more of it. B: unchanged if the added periods are plain mirror periods. TM:
same expectation, least certain.
\item \textbf{Value of a cancellation.} Removing a fraction $f$ gives $Q_{3\rm dB}/(1-f)$ and needs
$\ln[1/(1-f)]/(2\kappa\Lambda)$ more periods per side. From $Q=12\,903$: $f=0.1$: 14\,300 (+3 periods);
$f=0.3$: 18\,400 (+10); $f=0.5$: 25\,800 (+20). For scale, the TM comb pair reached $f\approx0.23$.
\end{itemize}

\section{Plan (nothing dispatched)}
\begin{enumerate}\setlength{\itemsep}{1pt}
\item Plain TE, $N=98$, monitors spanning the whole device (110\,$\mu$m), plus the complex field in
$\pm2$\,$\mu$m around the $\pi$ shift. Decides the end loss and measures the junction source.
\item Plain TE, $N=166$, surface fields stored. Pass: every harmonic group within 2 points of the $N=98$ table,
correlation above 0.95.
\item One perturbation of the innermost teeth chosen by Eq.~(\ref{eq:ls}); compare predicted and simulated
radiation change, resonance, width, $P_{\rm rad}/U$. Decides between the 83\,\% ceiling and the 10\,\% estimate.
\item Only if item 3 agrees: four to six junction knobs ($\pm h$ each, about ten short runs), solve, validate on
one unseen geometry, then one Q3dB device.
\item For B: give the inverse design the radiated power objective with the corrected projection.
\end{enumerate}
Calibration rules (v6.2): track the resonance, both mesh levels, measure the noise floor, control run with the
scatterer explicitly off; a higher peak $T$ alone is not evidence, report $Q_i$, $Q_c$, width and resonance
together.

\section{References}
Johnson, Fan, Mekis, Joannopoulos, APL 78, 3388 (2001): the mechanism. Johnson et al., PRE 65, 066611 (2002):
shifting boundaries. Srinivasan, Painter, Opt.\ Express 10, 670 (2002); Englund et al., Opt.\ Express 13, 5961
(2005); Quan, Lon\v car, Opt.\ Express 19, 18529 (2011): light cone design. Tran et al., PRB 79, 041101 (2009):
band folding. Lalanne et al., Laser Photonics Rev.\ 12, 1700113 (2018): the rigorous form of
Eq.~(\ref{eq:split}). Alaee et al., Opt.\ Commun.\ 407, 17 (2018): exact multipole moments. Full list,
derivations and numerical checks: v6.2; detailed evidence for this version:
\texttt{docs/radiation\_cancellation\_review\_2026-10-05.pdf}.
\end{document}
